\documentclass[10pt,a4paper]{amsart}
\usepackage[margin=19mm]{geometry}
\usepackage{amsmath,amssymb,mathtools}
\usepackage[scr=rsfs,cal=euler]{mathalfa}
\usepackage{xfrac}
\usepackage[unicode,hidelinks,bookmarks=false]{hyperref}
\newcommand{\ie}{i.e.\ }
\newcommand{\cf}{cf.\ }
\newcommand{\resp}{respectively\ }
\newcommand{\Subsection}{\subsection}
\providecommand{\nobreakdash}{\nobreak\hbox{-}}
\setlength{\emergencystretch}{2em}
\multlinegap0pt
\usepackage{bm}
\usepackage{bbm}
\usepackage{dsfont}
\usepackage{xspace}
\usepackage{cancel}
\usepackage{mathtools}
\usepackage{amsthm}
\makeatletter
\@ifundefined{theorem}{\newtheorem{theorem}{Theorem}}{}
\makeatother
\title[Small-2 Sets Are Riesz Sets]{Small-2 Sets Are Riesz Sets}
\author{A. To-Ming Lau, A. \"Ulger}
\date{Source: arXiv:2606.30570v2 (CC BY 4.0)}
\begin{document}
\maketitle

\noindent\emph{Source and licence.} Small-2 Sets Are Riesz Sets. Version arXiv:2606.30570v2 (CC BY 4.0), distributed under the Creative Commons Attribution 4.0 International licence, as recorded in the arXiv licence metadata for this exact version and shown on the abstract page https://arxiv.org/abs/2606.30570v2. Full rights notice: licenses/arxiv-2606.30570-CC-BY-4.0.txt.\par\smallskip

\noindent\textit{Editorial scope.} Counted unit: the Theorem whose source label is \texttt{None} in the article (source0.tex lines 222--223, source label \texttt{None}), delivered together with the article's own complete original proof of it (source0.tex lines 224--246, one \texttt{proof} environment of 23 lines). No mathematics is written, repaired or re-derived: statement and proof are the article's own text line for line. Required context retained uncounted: none. Every definition, display and label of those lines is retained unchanged. Omitted from this package and not redistributed: 1--221, 247--652. Nothing inside a retained excerpt is omitted or altered: every retained body is delivered exactly as the article's lines are written, so no omission receipt of any kind accompanies this package. Citations: the retained excerpts carry no citation command at all, so no bibliography entry of the article is transcribed and no reference is redistributed. Reference adaptation: none was needed and none was made; every reference inside the delivered unit resolves inside it, so no unresolved reference or citation is produced. Printed slips kept literal: no printed word, symbol, hypothesis or number of the counted statement or of its proof was altered. Layout and class: the article's own class is \texttt{article} with packages including amsmath, amsthm, amssymb; the wrapper is the lane's pinned \texttt{amsart} editorial wrapper at 10pt on A4, which restates the theorem-like environments the retained text needs and adds the article's own macro definitions that the retained text needs (none), with extra wrapper packages (bm, bbm, dsfont, xspace, cancel, mathtools). The delivered numbering is the unit's own. Assets: no retained span contains a figure environment, a graphics-inclusion command, a PDF-page inclusion, a table or a TikZ picture; no image file is copied into this package, the package declares no asset dependency and no image file is redistributed with it. Licence: version arXiv:2606.30570v2 of this article is distributed under the Creative Commons Attribution 4.0 International licence, as recorded in the arXiv licence metadata for this exact version and shown on the abstract page https://arxiv.org/abs/2606.30570v2. A full rights notice is delivered at licenses/arxiv-2606.30570-CC-BY-4.0.txt. Characters: every retained excerpt is pure ASCII, so the delivered engine, XeLaTeX with its default Latin Modern text font, reports no missing character.
\begin{theorem} Suppose that $ A^{**}A^{**}\subseteq A $. Then $ B^{**}B^{**}\subseteq A\subseteq B $ so that the algebra $ B^{**} $ is Arens regular and an ideal in its bidual.	
\end{theorem} 

\begin{proof}
	Before proving the inclusion $ B^{**}B^{**}\subseteq A $, let us first prove that the inclusion $ A^{**}A^{**}\subseteq A $ implies that the algebra $ B $ is Arens regular and an ideal in its bidual.\vspace{2mm}
	
	\noindent Let $ j:X\rightarrow A^{*} $ be the natural injection, $ j(f)=f $. As noted above, $ j $ is an isometry. Then $ j^{*}:A^{**}\rightarrow B $ is a surjective Banach algebra homomorphism that extends the natural injection $ A\hookrightarrow B $. That is, for $ a\in A $,  $ j^{*}(a)=a $. It follows that $ B $ is isomorphic to the quotient algebra $ A^{**}/Ker(j^{*}) $ of $ A^{**} $. So, since by the preceding theorem, the algebra $ A^{**} $ is Arens regular and an ideal in its bidual,  the algebra $ B $ also is Arens regular and an ideal in its bidual.\vspace{2mm}
	
	Next let us see that $ B^{**}B^{**}\subseteq A $. Since $ (B,X) $ is a dual algebra,\vspace{2mm}
	
	 $ B^{**}=B\oplus X^{\perp} $.\vspace{2mm}
	 
	 \noindent Let $ \tilde{u}=u+\lambda $ and $ \tilde{v}=v+\mu $ be two elements of $ B^{**} $ in this decomposition. Then, \vspace{3mm}
	
	$ \tilde{ u}\tilde{v}=uv+u\mu+v\lambda+\lambda \mu $.\vspace{3mm}
	
	\noindent Since $ (B,X) $ is a dual algebra,  for $ f\in X $, the functionals  $ u\cdot f$  and $v\cdot f$ are in $ X $. Since $ B $ is an ideal in $ B^{**} $, the products $ u\lambda $ and $ v\mu $ are in $ B $. Hence, for each $ f\in X $, $ <u\lambda,f>=<\lambda,u\cdot f>=0 $. So $ u\lambda=0$. For the same reasons, $v\mu=0 $ too. To see that $ \lambda\mu=0 $ too, let $ (u_{i}) $ be a net in $ B $ that converges to $ \lambda $ in the weak-star topology of $ B^{**} $. Since, as just seen, $ u_{i}\mu=0 $, we conclude that $ \lambda\mu=0 $ too. So\vspace{2mm}
	
	$ \tilde{u}\tilde{v}=uv$.\vspace{2mm}
	
	\noindent Hence $ B^{**}B^{**}\subseteq BB $.\vspace{ 2mm}
	
	To see that $ BB\subseteq A $, let $ u,v $ be two elements in $ B $ with $ ||u||\leq 1 $. Let $ (a_{i}) $ be a net in $ A_{1} $ that converges in the weak-star topology of $ B $ to $ u $. Then, since $ B $ is an ideal in its bidual, the multiplication operator $ L_{v} $ is weakly compact on $ B $, so on $ A $. Then the net $ (a_{i}v) $, which is contained in $ A $, is relatively weakly compact. So it converges weakly to $ uv $, so that $ uv $ is in $ A $. Thus,\vspace{2mm}
	
	$ B^{**}B^{**}\subseteq BB\subseteq A $.
	\end{proof}

\end{document}
